% =====================================================================
%  sec03 -- Obstruction two: Delta_dyn > 0 does not imply m > 0
% =====================================================================
\section[Obstruction two: the dynamical block does not imply a gap]{Obstruction two: $\Dyn>0$ does not imply $m>0$}\label{sec:obs2}

\subsection{The definitions have to be pinned first, or the comparison is meaningless}

Two quantities are being compared, and comparing them without pinning both is what produced the
confusion the corpus is now unwinding. Let
\begin{equation}
\Gpp:=-\partial_a^{2}\log\rho(a)\big|_{a=0},
\label{eq:gdef5}
\end{equation}
where $\rho$ is the marginal law of the coarse variable $a$ under the model's own measure, so that
in zero dimensions $\rho\propto e^{-S}$ gives $\Gpp=S''(0)$, and in $0+1$ dimensions
$\rho=|\psi_0|^{2}$ gives the ground-state marginal. This is the same $-\partial^{2}\log\rho$ that
the conditional-expectation expansion of \eqref{eq:split5} computes, and the two agree numerically
to $1.1\times10^{-16}$ on the single plaquette. \LT{} For the gap, let
\begin{equation}
m:=\inf\bigl(\mathrm{spec}(H)\setminus\{E_0\}\bigr)-E_0 ,
\label{eq:mdef5}
\end{equation}
\emph{which is defined only when the model has a transfer direction}. \LT{} In a
zero-dimensional model --- a single plaquette, a $0+0$ double well --- there is no transfer
matrix and no $E_1$, so $m$ is \emph{undefined}: not $0$, not $>0$. This is a discipline and not a
pedanticity, because both misreadings fabricate a result. Writing ``undefined $=0$'' manufactures
a counterexample to the implication out of a model that says nothing; writing ``undefined $>0$''
defends the implication for free. The corpus records this explicitly, and it is the reason model
M1 below is scored \emph{undecidable} rather than either way.

\src{\code{T03-TOY.md} \S1 (both definitions, verbatim, including the two-sentence discipline on
``undefined''); \code{Q11-CERT.md} \S2 (the anchor $S-D=+0.135994711\ldots$ reproduced to
$1.1\times10^{-16}$ by the $\Gpp=-\partial^{2}\log\rho$ route).}

\subsection{Six models, one table}

The following is the whole comparison. It is deliberately small: three of the six are vacuous,
and the two that matter are exact.

\begin{center}
\footnotesize
\begin{tabularx}{\textwidth}{@{}L{0.55cm}L{4.6cm}L{2.45cm}L{2.5cm}L{2.4cm}@{}}
\toprule
\textbf{\#} & \textbf{Model} & $\Gpp$ & $m$ & \textbf{Verdict} \\
\midrule
M1 & one plaquette $\mathrm{SU}(2)$, $\beta=2.5$, $\theta=\pi$ & $+0.135994711$ & undefined
   (no transfer direction) & vacuous \\
M2 & 2D Yang--Mills, Migdal exact solution & $0$ \emph{exactly} & $3g^{2}/8=0.375>0$
   & \textbf{reverse fails} \\
M3 & 1D Dirichlet string, lattice $N=999$ & $\pi^{2}/L^{2}>0$ & $\pi/L>0$ & consistent,
   but a Gaussian identity \\
M4 & flat double well $a^{4}/4-a^{2}/2$ & $0+0$: $-1$ exact; $0+1$: $+0.5889$ &
   $0+0$: undefined; $0+1$: $0.7250$ & $0+1$ consistent \\
M5 & free field $S\equiv0$ & $0$ & $0$ & vacuous \\
\textbf{M6} & \textbf{harmonic $\otimes$ free}
   $\bigl(p_a^{2}+p_b^{2}\bigr)/2+\omega^{2}a^{2}/2$ & $\mathbf{2\omega=2>0}$ & $\mathbf{0}$
   & \textbf{counterexample} \\
\bottomrule
\end{tabularx}
\end{center}

\src{\code{T03-TOY.md} \S2, table (all six rows, with the same entries); \code{LOOP-RUNBOOK.md}
\S12.3 (the same table in the campaign log).}

\subsection{M6: the counterexample, and why it is not a cheap special case}

In M6 the $a$-direction is a harmonic oscillator and the $b$-direction is free. The marginal is
$\rho_a\propto e^{-\omega a^{2}}$, so $\Gpp=2\omega=2>0$ identically; the spectrum is
$\mathrm{spec}(H)=[\omega/2,\infty)$, continuous from $E_0$ upward, so $m=0$. \LT{} Both statements
are exact --- the first is a Gaussian integral, the second an exact spectrum --- and neither is a
numerical observation.

The reason this is the \emph{minimal witness} rather than a trick is structural, and it is the
sentence that should be quoted whenever the route is proposed again: $\Gpp$ is the curvature of
\emph{one} direction's marginal, while $m$ is a property of the \emph{whole} spectrum. Any model
whose marginal has a local maximum at the symmetric point while possessing one flat (massless)
direction will produce $\Gpp>0$ together with $m=0$. There are infinitely many such models, and
they are not contrived: the mismatch is one structural fact, not one exotic example.

\src{\code{T03-TOY.md} \S2 (M6), \S2 ``why this is not a cheap special case'' (verbatim);
\code{LOOP-RUNBOOK.md} \S12.1 (the same counterexample recorded as the campaign-level verdict).}

\subsection{The mechanism, which is more useful than the counterexample}

The counterexample is an instance of an identity. In $0+1$ dimensions with a flat measure,
$H=p^{2}/2+V$, $\rho=|\psi_0|^{2}$, so
\begin{equation}
\Gpp=-2\,\frac{\psi_0''}{\psi_0}\Big|_{0}=4\bigl(E_0-V(0)\bigr),
\label{eq:mech5}
\end{equation}
by inserting the Schr\"odinger equation $-\psi_0''/2+V\psi_0=E_0\psi_0$. \LT{} Two numerical
checks are recorded: the oscillator gives $4(\omega/2-0)=2\omega$ against the measured
$2.000006$; the double well gives $0.5889400$ by differencing $\psi_0$ against
$4(E_0-V(0))=0.5889397$, a relative agreement of $3.9\times10^{-7}$ (grid error, $h=3\times10^{-3}$).

The reading of \eqref{eq:mech5} is the whole of \S\ref{sec:obs2}: $\Gpp>0$ is \emph{equivalent} to
$V(0)<E_0$, i.e.\ to the statement that the point $a=0$ lies above the ground-state energy --- a
classically forbidden point, a barrier top, a local maximum of the marginal density. It is a
statement about \emph{one point's value} against \emph{one number}. The gap $m=E_1-E_0$ is
determined by the \emph{global} shape of the potential and the boundary conditions. To infer
$m>0$ from $\Gpp>0$ is to claim that a single-point potential value exceeding the ground-state
energy implies a gap above it, and \eqref{eq:mech5} plus M2/M6 show that this is false in both
directions.

\src{\code{T03-TOY.md} \S3 (the proposition and its two-line proof; both numerical checks;
the ``two independent objects, two different propositions'' reading, verbatim).}

\subsection{The reverse direction fails too, and exactly}

The implication is also false backwards. In two-dimensional Yang--Mills on a spatial circle the
Migdal summation reduces the transfer operator to quantum mechanics on the group manifold,
$H=(g^{2}/2)C_2$ with $E_R=(g^{2}/2)C_2(R)$; the ground state is the trivial representation, i.e.\
the constant function with respect to Haar measure, so \emph{every} marginal of the coarse
variable is flat, $\partial^{2}\log\rho\equiv0$, and $\Gpp=0$ exactly. Meanwhile
$m=(g^{2}/2)C_2(\text{fund})=g^{2}(N^{2}-1)/(4N)$, which at $N=2$, $g^{2}=1$ is $3g^{2}/8=0.375>0$.
\LT{}$^{*}$ \,\, The asterisk is honest bookkeeping: the model derivation is strict, but the
statement $H=(g^{2}/2)C_2$ is quoted as the standard 2D result and its literature source was not
checked verbatim, so the label is \LT{} inside the model and \LC{} as a citation. The reverse
failure does not depend on this row: M6 already establishes the structural decoupling of $\Gpp$
from $m$ independently.

\src{\code{T03-TOY.md} \S2 (M2), \S5 (``\LT{}$^{*}$: the model derivation is strict, the literature
citation was not checked verbatim'').}

\subsection{The four routes from the curvature to a gap: all four are broken}

It is not that one step is missing. The implication has four candidate proofs in the corpus, and
each fails for a specific reason; they are recorded here because each failure is a distinct
theorem.

\begin{enumerate}
\item \textbf{P1: $\Gpp>0\Rightarrow\fone>0\Rightarrow$ atom $\Rightarrow$ isolation
$\Rightarrow$ exponential $\Rightarrow$ gap.} \LN{} --- not ``unproved'' but \emph{reversed}. The
counterexample is exact and is 2D Yang--Mills again: it has an exact area law
$\sigma_R\propto C_2(R)>0$ and a gapped discrete spectrum, hence $H1$ holds, hence $\Gpp>0$ by
the identity, and yet $\fone(\tr F^{2})=0$ exactly, because $\tr F^{2}=\tr E^{2}=H$ is conserved
and the characters are orthogonal. The direction is also wrong in strength: $\fone>0$ together
with isolation is strictly stronger than a gap, so using the weaker hypothesis to reach the
stronger conclusion upgrades the conclusion into the premise. What the route needs is the
existence of a sub-threshold atom, $\Lstar>m_1^{2}$, which is condition I2 --- one level
\emph{above} Clay, not below it (\S\ref{sec:target}).
\item \textbf{P2: $\Gpp>0\Rightarrow$ some correlator does not decay.} \LT{} inside a channel and
open across channels. Within one channel, K\"all\'en--Lehmann makes
$m^2_{\mathrm{eff}}(k\to0)\ge\lambda$ and $\langle OO\rangle_c\le Ce^{-\lambda|x|}$ the same
sentence in two languages, so the step carries zero information. Across channels it is a
different statement entirely: $\Gpp$ is conditional on the coarse amplitude $a$ and lives in the
$S_W$-channel, while Clay needs the $0^{++}$ channel of $\tr F^{2}$; the kinematic bridge is
excluded by $Z_N^{(1)}$ charge grading, and the only dynamical bridge the corpus contains is the
$u$-identity. The missing sentence --- ``positive curvature in the coarse channel implies that the
spectral measure of the $\tr F^{2}$ channel has support above zero'' --- is not written anywhere in
the corpus. It is the sector-transfer problem, and it has a name: $e_3$.
\item \textbf{P3: $\Gpp>0\Rightarrow$ transfer-matrix spectral gap.} \LT{} as a renaming and \LN{}
as an argument. The renaming is the four-name equivalence of \S\ref{sec:obs1}. The circularity is
that the only remaining content, the uniformity $\inf_V\Delta_V>0$ over volumes, is \emph{literally
the gap}: for finite $V$, $\Delta_V>0$ is trivial by Perron--Frobenius/Jentzsch, and the
infimum over $V$ being positive is the statement to be proved.
\item \textbf{P4: $\Gpp>0\Rightarrow\langle\tfrac14\tr F^{2}\rangle>0\Rightarrow\cdots$} --- no
chain. The second half has never been written, and the first half is refuted in 2D (where
$\langle H\rangle>0$ holds for a gapped theory while $\fone(\tr F^{2})=0$). In $4$D the condensate
is an operator-product-expansion quantity requiring subtractions, i.e.\ \LC{} at best, and there
is no proved channel between it and an atom of $\nu$. The corpus's phrasing ``the same object
under several names'' is a programme statement, not a theorem.
\end{enumerate}

\src{\code{T01-IMPLY.md} \S2 (the four-row table of what each route additionally needs), \S3
(the four verdicts with their sources), \S3 P1 (the 2D counterexample and the strength reversal);
\code{IMP-ARROWS.md}:33--35 (the thin-atom counterexample, machine-checked as
\code{H1_does_not_imply_atom}), \code{:57} (row 5(iv)), \code{:60} (row 8, K\"all\'en--Lehmann both
ways), \code{:76} ($Z_N^{(1)}$ grading), \code{:92--95} (the 2D exact counterexample),
\code{:113} ($e_3$), \code{:137--145} (chain III and $k_2$); \code{ch06-tools.tex}:31--48
(multiplicative residue transport), \code{:74--75} ($\langle\tfrac14\tr F^{2}\rangle$ as an OPE
quantity).}

\subsection{The domain question, answered honestly}

A careful reader will object that within the identity's domain the two objects are the \emph{same
number}: $\Gpp=1/\Var(a)=m^2_{\mathrm{eff}}(k\to0)$ is model-independent algebra, so no
counterexample can exist inside pure Yang--Mills at fixed $(L,a)$. That objection is correct and
is why the obstruction is stated in the form it is. Two things follow, and both are needed.

\emph{First}, the bare implication is nevertheless false in the natural class in which the route
was being pursued. Take $G_{\mathrm{YM}}\times U(1)_{\mathrm{weak}}$ with the abelian factor in the
Coulomb phase and the coupling weak enough: the Yang--Mills factor has an area law, reflection
positivity holds, and its $a$-channel curvature is positive; but the full theory's spectrum has no
gap, because of the photon sector. Hence ``$\Gpp>0$ together with $m=0$'' occurs in the class
``reflection-positive lattice gauge theory''. \LN{} \emph{Second}, for pure Yang--Mills the failure
therefore cannot be a sign failure at fixed volume; it must be the failure of \emph{uniformity} ---
finite volume gapped, infinite volume not --- which is precisely where \S\ref{sec:obs1} put the
break. Both routes to a rescue end at the same wall: the object must carry its quantifiers, i.e.\
the statement must be about a uniform positive lower bound plus purity (no decoupled neutral
massless sector). That conjunction is again $e_3$.

\src{\code{T01-IMPLY.md} \S4.3 (the $G_{\mathrm{YM}}\times U(1)_{\mathrm{weak}}$ class
counterexample, and the ``(c) does not hold, (b) does'' verdict); \code{IMP-ARROWS.md}:83--90
(\S2.1(a$'$), the same construction).}

\subsection{What this section establishes}

\begin{itemize}
\item \LT{} $\Gpp$ and $m$ are decoupled in both directions: M6 gives $\Gpp=2\omega>0$ with $m=0$;
      2D Yang--Mills gives $\Gpp=0$ with $m=3g^{2}/8>0$.
\item \LT{} The mechanism is the identity $\Gpp=4(E_0-V(0))$: the curvature measures whether
      $a=0$ lies above the ground-state energy, a single-point comparison with no functional
      relation to the spectrum.
\item \LN{} All four routes from the curvature to a gap are closed: P1 is reversed (2D exact),
      P2 is a renaming inside one channel and open across channels ($e_3$), P3 is circular
      ($\inf_V\Delta_V>0$ is the gap), P4 has no chain.
\item \LM{} For pure Yang--Mills the surviving difficulty is not the sign of $\Gpp$ but its
      \emph{uniformity} in $(L,a)$ jointly with purity; that conjunction is the exact missing
      step, and it is not a step on this route at all.
\end{itemize}
